\chapter*{Vorwort}
\addcontentsline{toc}{chapter}{Vorwort}

\begin{merksatz}
\boxtitel{Worum es geht}
Ein Herz schlägt siebzig Mal in der Minute, und niemand zählt vor. Eine Nervenzelle feuert, wenn der Reiz eine Schwelle überschreitet, und schweigt darunter. Eine Epidemie wächst, kippt und verebbt. Ein Wetter, das man heute kennt, ist in zwei Wochen nicht mehr zu berechnen --- nicht weil die Gleichungen fehlen, sondern obwohl man sie hat.

\begin{center}
\begin{minipage}{0.9\textwidth}
\textbf{Dieser Band handelt von Systemen, deren Zukunft aus ihrer Gegenwart folgt --- und davon, wie weit diese Folge trägt.}
\end{minipage}
\end{center}

\begin{center}
\begin{minipage}{0.88\textwidth}
\kursiv{Band 14 hat Differentialgleichungen gelöst; Band 18 hat mit ihnen Populationen beschrieben und am Ende die logistische Abbildung gesehen, die ins Chaos führt. Dieser Band nimmt den Faden dort auf und fragt nicht mehr nach der Lösung, sondern nach der \Begriff{Gestalt} der Lösungen: Wohin laufen sie, wenn man lange genug wartet? Was geschieht, wenn ein Parameter sich verschiebt? Wann wird aus Ruhe eine Schwingung, aus Schwingung ein Durcheinander --- und ist das Durcheinander Zufall oder Gesetz?}
\end{minipage}
\end{center}
\end{merksatz}

\begin{merksatz}
\boxtitel{Der Leitfall: das Herz}
Eine Herzmuskelzelle ruht bei $-85$ Millivolt. Ein Reiz, der eine Schwelle überschreitet, löst ein Aktionspotential aus --- einen Ausschlag auf $+20$ Millivolt, der etwa $300$ Millisekunden dauert und sich selbst beendet. Der Sinusknoten erzeugt solche Reize von allein, siebzig Mal in der Minute. Das Reizleitungssystem trägt sie durch das Herz, und Millionen Zellen antworten im Takt.

\begin{center}
{\small
\renewcommand{\arraystretch}{1.5}
\begin{tabular}{@{}P{3.0cm}P{5.8cm}@{}}
\toprule
\textbf{Die Zelle} & Ein erregbares System: Ruhe, Schwelle, Antwort, Rückkehr --- Kapitel~1 bis~3.\\
\textbf{Der Taktgeber} & Ein Grenzzyklus, der von allein entsteht, wenn ein Parameter kippt --- Kapitel~4 und~5.\\
\textbf{Der Schrittmacher} & Ein Oszillator, dem von außen ein Takt aufgezwungen wird --- Kapitel~6.\\
\textbf{Das Stolpern} & Alternans: Wenn jeder zweite Schlag anders ist, verdoppelt sich die Periode --- Kapitel~7.\\
\textbf{Das Flimmern} & Wellen, die sich in Spiralen zerlegen, bis kein Takt mehr da ist --- Kapitel~13.\\
\textbf{Die Messung} & Was man aus einer Herzfrequenzkurve über das System lernen kann --- Kapitel~10 bis~12.\\
\bottomrule
\end{tabular}}
\end{center}

\begin{center}
\begin{minipage}{0.88\textwidth}
\kursiv{\textbf{Das Herz ist der Leitfall, weil es alle Erscheinungen dieses Bandes auf engem Raum zeigt} --- und weil jede davon klinisch zählt: Die Schwelle ist die Grundlage der Defibrillation, der Grenzzyklus der des Schrittmachers, die Periodenverdopplung ein Vorbote des plötzlichen Herztods, die Spiralwelle der Mechanismus des Kammerflimmerns. Die Gleichungen sind einfach genug, um sie vollständig zu verstehen, und reich genug, um nicht vorhersagbar zu sein.}
\end{minipage}
\end{center}
\end{merksatz}

\begin{merksatz}
\boxtitel{Drei Ergebnisse, die man nicht erwartet}
\begin{center}
{\small
\renewcommand{\arraystretch}{1.55}
\begin{tabular}{@{}P{1.0cm}P{7.4cm}@{}}
\toprule
$1$ & Drei Gleichungen mit drei Unbekannten, ohne jeden Zufall, erzeugen einen Verlauf, der sich nie wiederholt und nach wenigen Zeiteinheiten nicht mehr vorhersagbar ist. \kursiv{(Kapitel 8)}\\
$2$ & Der Weg von der regelmäßigen Schwingung ins Chaos führt über Periodenverdopplungen, deren Abstände sich stets um denselben Faktor $4{,}669\ldots$ verkürzen --- bei jedem System, in dem er vorkommt. \kursiv{(Kapitel 7)}\\
$3$ & Aus einer einzigen Messgröße --- der Herzfrequenz, Schlag für Schlag --- lässt sich die Gestalt des ganzen Systems zurückgewinnen, das sie erzeugt hat. \kursiv{(Kapitel 10)}\\
\bottomrule
\end{tabular}}
\end{center}

\begin{center}
\begin{minipage}{0.88\textwidth}
\kursiv{Das erste Ergebnis ist der Grund, warum es diesen Band gibt. \textbf{Determinismus und Vorhersagbarkeit sind zwei verschiedene Dinge}, und der Unterschied ist keine Frage der Rechenleistung, sondern der Gestalt des Systems. Band 27 hat gezeigt, was der Rechner vom Zufall weiß; dieser Band zeigt, wo der Rechner --- und jeder andere --- mit dem Wissen am Ende ist.}
\end{minipage}
\end{center}
\end{merksatz}

\begin{merksatz}
\boxtitel{Die sechs Teile}
\begin{center}
{\small
\renewcommand{\arraystretch}{1.5}
\begin{tabular}{@{}P{1.4cm}P{3.6cm}P{3.0cm}@{}}
\toprule
\textbf{Teil} & \textbf{Thema} & \textbf{Kapitel}\\
\midrule
I & Der Fluss & $1$--$3$\\
II & Wenn Parameter wandern & $4$--$6$\\
III & Der Weg ins Chaos & $7$--$9$\\
IV & Erkennen und Messen & $10$--$12$\\
V & Raum und Zeit & $13$--$14$\\
VI & Stand & $15$\\
\bottomrule
\end{tabular}}
\end{center}

\schritt{Was vorausgesetzt wird}
Differentialgleichungen und ihre Lösung (Band 14), lineare Algebra mit Eigenwerten (Band 2), die logistische Abbildung und die Lotka-Volterra-Gleichungen (Band 18), Wahrscheinlichkeit und Zeitreihen in dem Maß, wie Band 25 und 26 sie eingeführt haben. Alles andere wird hier entwickelt, und wo ein Beweis übergangen wird, steht es dabei --- die Auslassungen sind im Anhang gesammelt.

\begin{center}
\begin{minipage}{0.88\textwidth}
\kursiv{\textbf{Jedes Kapitel endet mit einem Bruchfall}: einem Beispiel, an dem die Methode des Kapitels versagt --- die Linearisierung am Zentrum, die Verzweigung bei fehlender Symmetrie, die Einbettung bei farbigem Rauschen. Und jedes Kapitel hat drei Übungsaufgaben mit vollständiger Lösung und Probe; die Zahlen darin stammen aus Rechnungen, die der Leser mit Band 27 nachvollziehen kann.}
\end{minipage}
\end{center}
\end{merksatz}
